\subsection{Metric definitions and aggregation}\label{sec:metric_protocol}
The five-sequence evaluation distinguishes mathematical optima, achieved anchor displacement, measured raster errors and geometric motion outcomes. Each quantity is aggregated over its stated frame, feature or event population.

\paragraph{Normalized Reference-Order Error (NROE, $P$).}
For each frame, stable sorting of displayed anchors gives $\pi$. The point-model optimum is $\tau_{\rm pt}(\pi)=\frac12\max\{0,\max_{r<s}(q_{\pi_r}-q_{\pi_s})\}$. Its temporal mean divided by the fixed sequence reference extent $L$ is $P$. This removes widths, gaps, canvas bounds, and anchor radii, permitting a monotone calibration of distinct anchors. Stable tie-breaking assigns separate calibrated values to tied pixels; it is an optimistic diagnostic in that case. $\tau^*_{\rm pt}$ minimizes this point cost over tree-consistent orders. Neither $P$ nor $\tau^*_{\rm pt}$ equals achieved reference displacement under the full interval model.

\paragraph{Achieved continuous displacement.}
For each retained feature-frame anchor define $e_i^t=|u_i^t-q_i^t|$. Report
\[
 E_{\rm pool}=\frac{\sum_t\sum_{i=1}^{n_t}e_i^t}{L\sum_tn_t},\qquad
 E_\infty=\frac{\max_{t,i}e_i^t}{L}.
\]
The original record also saves the pooled 95th percentile. Pooling weights frames by their retained feature count; the maximum weights no frame average. The full certificate $\tau_t^*$ instead bounds a frame's maximum, so it must not be compared directly with $E_{\rm pool}$. Values here use continuous RG-MTM-P and RG-MTM-R anchors in reference units, before raster rounding. TMTM/ST-MTM native axes have different units and width models and are not assigned these direct-reference scores without an explicit calibration protocol.
The achieved displacement values are reported in main Table 2; the accompanying numerical ledger also retains the pooled 95th percentile.

\paragraph{Normalized Merge-Level Error (NMLE, $D$).}
$\Delta_F=\max_{t,x}F_t(x)-\min_{t,x}F_t(x)$ is fixed per processed sequence. $D_t=d_{\rm top}(\psi_t,T_t)/\Delta_F$ measures the maximum pairwise labelled merge discrepancy in frame $t$. $D_{\rm mean}=K^{-1}\sum_tD_t$ averages frames; $D_{\rm max}=\max_tD_t$ is the worst frame. These are displayed as NMLE$_{\rm mean}$ and NMLE$_{\rm max}$; the symbols and saved numerical keys are unchanged. The fixed-label cophenetic basis follows existing labelled-tree metrics~\cite{Munch:2019:LINF}; the experiment-specific normalization does not create a new unlabelled tree distance. One retained leaf gives zero pairwise error. For split trees the profile and original levels are negated before taking maxima. Native-curve barriers are measured on unchanged sample traversals. The revised main table uses \path{results/filling.json}'s measured optimal-filling summaries, while \path{results/eval_projection.json}'s \texttt{dstar} is the fixed-order closed-form optimum. Their small clearance/raster discrepancies are checked separately rather than silently identifying an optimum with a measurement.

\paragraph{Motion events and inference.}
A pair of tracks must be present at $t$ and $t+k$; its true distance is the Euclidean distance between retained extrema. A clear event has distance change exceeding $\epsilon_{2D}=0.02\,\mathrm{diag}(\Omega)$ in magnitude. Main results fix $k=2$. Displayed pair distance uses actual anchors and a method-specific threshold $\epsilon_m=0.02A_m$, where $A_m$ is that method's full-sequence occupied anchor range. A nonzero opposite sign is a reversal; a change within $[-\epsilon_m,\epsilon_m]$ is a miss. ME is $(N_{\rm rev}+N_{\rm miss})/N_{\rm clear}$. The same true events are compared across methods, but the no-change thresholds depend on declared display ranges. If $N_{\rm clear}=0$, error rates are undefined; they are not perfect scores. Motion labels are computational geometry, not participant responses. Projected-reference events use a different truth and threshold and must not be subtracted from 2D-truth rates to estimate projection loss.

Paired differences reuse the same events and employ circular moving-block bootstrap over time, with block length $k+2=4$ at the primary setting. The retained intervals are exploratory and have no multiplicity correction. The historical sparse January PM2.5 events cannot establish a seasonal superiority claim; threshold/window and fixed-direction results remain controls.

\paragraph{Geometry and stability diagnostics.}
For unordered extremum pairs in a frame, set $a_{ij}=|u_i-u_j|$, $b_{ij}=\|R_i-R_j\|$, and $\alpha=(\sum a_{ij}b_{ij})/(\sum a_{ij}^2)$. The implementation averages
\[
 S_t=\sqrt{\frac{\sum_{i<j}(\alpha a_{ij}-b_{ij})^2}{\sum_{i<j}b_{ij}^2}}
\]
over frames with at least three leaves. NN is the framewise mean agreement of each extremum's nearest neighbour under $a$ and $b$, then averaged over the same frames. Stable array order resolves distance ties as implemented. Report the number of eligible frames; degenerate cases need explicit treatment if this protocol is extended beyond the present suite.

Flips counts sign changes in display order among pairs matched across consecutive frames, divided by the number of eligible matched pairs. $J$ averages $|u_i^t-u_j^{t-1}|/A_m$ over individual matches. These stability diagnostics do not compare motion with a source target: true crossings or rapid source motion can increase either score. The QP's distance objective uses support centroids; the retained diagnostic $S$ uses extremum locations. These choices answer different questions and are not silently interchanged.

ST-MTM's stated SNS has squared residuals and centroid distances, TW penalizes introduced neighbours by original-space ranks, and TD sums matched displacements. SNS and TW are adopted from prior work~\cite{Smelser:2025:Scale,Venna:2001:NP}, while ST-MTM specifies TD for its matched display positions. NN is not TW, $J$ is not TD, and $(\mathrm{mean}\,S_t)^2$ is not $\mathrm{mean}\,(S_t^2)$. The separate standard-metric re-evaluation below uses the required entities and aggregation without changing these historical diagnostics.

\paragraph{Source mapping.}
Order statistics use \path{results/final_performance.json}; motion and geometry use \path{results/eval_projection.json}; measured barrier costs use \path{results/filling.json}; achieved anchors use \path{results/replicate_<dataset>.json}'s \texttt{tradeoff} entries at caps 0 and 0.2, normalized by its \texttt{certificate.axis}. The main-paper ledger stores source hashes and unrounded values. No metric selection or bold-face ranking is inferred from a combined score.

\paragraph{Method-name mapping.}
RG-MTM means Reference-Guided Merge Tree Maps. RG-MTM-P is the hierarchy-preserving mode; RG-MTM-R is the relaxed mode. In the saved records, \texttt{A} denotes RG-MTM-P and \texttt{R 20\%} denotes RG-MTM-R at $\kappa=0.2\Delta_F$. Keys \texttt{R 2\%} and \texttt{R inf} identify the other eligibility settings. The current code directory is \texttt{RG-MTM/code}; historical \texttt{CD-MTM} labels identify the same relaxed method.
